\documentclass{beamer}

\usepackage{tikz}
\usepackage{pgfplots}

\title{Eigen-\textit{stuff}}
\author{Ahmed Mohammed}
\date{PLUS Session, June 2023}
\institute{Georgia Institute of Technology}

\usetheme{Berlin}
\usefonttheme{serif}
\usecolortheme{seagull}

\AtBeginSection[] {
	\begin{frame}
		\frametitle{Table of Contents}
		\tableofcontents
	\end{frame}
}

\DeclareMathOperator{\nullspace}{Nul}

\begin{document}
	\begin{frame}
		\maketitle
	\end{frame}

	\section{Opening: Introduction to Eigen-\textit{stuff}}
	\begin{frame}
		\frametitle{Main Idea}
		If \(A \in \mathbb{R}^{n \times n}, \vec{v} \in \mathbb{R}^n,\) and,
		\[A\vec{v} = \lambda \vec{v},\ \text{with} \ \vec{v} \ne \vec{0}, \ \lambda \in \mathbb{R}\] then \(\vec{v}\) is an \emph{eigenvector} of \(A\) with a
		corresponding \emph{eigenvalue} of \(\lambda\).
	\end{frame}

	\begin{frame}
		\frametitle{A Picture (Before)}
		\begin{figure}[!h]
			\begin{tikzpicture}
				\begin{axis}
					[axis lines=middle, height=7cm, width=7cm, xmin=-10, xmax=10, ymax=10, ymin=-10,
					 xlabel=\(x_1\), ylabel=\(x_2\), grid=both, minor tick num=4, xticklabels={}, yticklabels={}, 
					legend style={nodes={scale=0.5, transform shape}}] % Thanks tex.stackexchange

					\addplot[->, domain=0:2, color=red!60!black] {x};
					\addlegendentry{\(\vec{v}_1\)}
					
					\addplot[<-, domain=-2:0, color=blue!60!black] {-x};
					\addlegendentry{{\(\vec{v}_2\)}}

					\addplot[->, domain=0:2, color=green!60!black] {2*x};
					\addlegendentry{\(\vec{v}_3\)}


					\addplot[<->, domain=-6:6, color=blue, dashed] {-x};
					\addplot[<->, domain=-4:4, color=green, dashed] {2*x};
					\addplot[<->, domain=-6:6, color=red, dashed] {x};
				\end{axis}
			\end{tikzpicture}
		\end{figure}
	\end{frame}

	\begin{frame}
		\frametitle{A Picture (After)}
		\begin{figure}[!h]
			\begin{tikzpicture}
				\begin{axis}
					[axis lines=middle, height=7cm, width=7cm, xmin=-10, xmax=10, ymax=10, ymin=-10,
					 xlabel=\(x_1\), ylabel=\(x_2\), grid=both, minor tick num=4, xticklabels={}, yticklabels={}, 
					legend style={nodes={scale=0.5, transform shape}}] % Thanks tex.stackexchange

					\addplot[->, domain=0:2, color=red!60!black] {x};
					\addlegendentry{\(A\vec{v}_1\)}
					
					\addplot[<-, domain=-1:0, color=blue!60!black] {-x};
					\addlegendentry{{\(A\vec{v}_2\)}}

					\addplot[->, domain=0:2.5, color=green!60!black] {(3.5/2.5)*x};
					\addlegendentry{\(A\vec{v}_3\)}

					\addplot[<->, domain=-6:6, color=blue, dashed] {-x};
					\addplot[<->, domain=-4:4, color=green, dashed] {2*x};
					\addplot[<->, domain=-6:6, color=red, dashed] {x};
				\end{axis}
			\end{tikzpicture}
		\end{figure}
	\end{frame}

	\begin{frame}
		\frametitle{Some Terminology}
		\begin{itemize}
			\item<1-> Each eigenvalue of a matrix has an infinite number of corresponding eigenvectors. This infinte amount of eigenvectors form an 
				\emph{eigenspace} of the matrix corresponding to \(\lambda\). \pause 
				(\(\text{Eigenspace of A Corresponding to } \lambda = \nullspace(A - \lambda I)\)).
			\item<2-> A basis for the eigenspace is the set of vectors found when finding the parametric vector form of \(\nullspace(A - \lambda I)\).
			\item<3-| alert@3> An \textbf{eigenbasis} is a set of \emph{eigenvectors} of a single matrix that form a basis for the space 
				(e.g., \(\mathbb{R}^n\)). Do not confuse this with the basis for the eigenspace!
		\end{itemize}
	\end{frame}

	\section{Main Activity: Some Questions}
	\begin{frame}
		\frametitle{Questions}
		\begin{itemize}
			\item<1-> Find the basis of the eigenspace of the following matrix corresponding to \(\lambda = 4\): 
					\[\begin{bmatrix}2.5 & 1.5 \\ 1.5 & 2.5\end{bmatrix}\]
			\item<2-> Find the other eigenvalue of the matrix, and find the basis for the corresponding eigenspace.
			\item<3-> How did I know that the matrix had two eigenvalues?
			\end{itemize}
	\end{frame}

	\begin{frame}
		\frametitle{Fall 2022, Midterm Exam \#2; Problem 9}
		Consider a Markov chain \(\vec{x}_{k + 1} = Px_{k}, \ k = 0, 1, 2, \dots\) with eigenvalues \(\lambda_1 = 1\), \(\lambda_2 = 1/2\), \(\lambda_3 = 0\).
		Let \(\vec{v}_1, \vec{v}_2\) and \(\vec{v}_3\) be eigenvectors corresponding to \(\lambda_1, \lambda_2, \lambda_3\), respectively:
		\[\vec{v}_1 = \begin{bmatrix}1 \\ 1 \\ 0\end{bmatrix}, \quad
		  \vec{v}_2 \begin{bmatrix}0 \\ -1 \\ 1\end{bmatrix}, \quad
		  \vec{v}_3 = \begin{bmatrix}-1 \\ 1 \\ 0\end{bmatrix}\]\pause
		\begin{itemize}
			\item<2->If \(\vec{x}_0 = \begin{bmatrix}1/4 \\ 1/2 \\ 1/4\end{bmatrix}\), then what is \(\vec{x}_1\)?
			\item<3->If \(\vec{x}_0\) is the same as before, then what is \(\vec{x}_k\)?
			\item<4->If \(\vec{x}_0\) is the same as before, then what is \(\vec{x}_k\) as \(k \to \infty\)?		
		\end{itemize}
	\end{frame}

	\section{Closing: Exam Review}
	\begin{frame}
		\frametitle{Check Understanding}
	Answer \textit{true} or \textit{false} to the following question:
		\begin{itemize}
			\item<1-> If the eigenspace of a stochastic matrix corresponding to \(\lambda = 1\) has a dimension of 1, then the corresponding Markov chain
					always converges to the same vector regardles of the initial state.
		\end{itemize} \pause 
		Complete Worksheet 2; work on Worksheet 3. 
	\end{frame}
\end{document}